\section{CONCLUSÕES E TRABALHOS FUTUROS}
\label{sec:conclusao}
O presente trabalho teve como objetivo central conceber e validar uma
arquitetura de \textit{pipeline} de ingestão de dados orientada a metadados,
voltada para a extração, transformação e carga de arquivos fiscais,
especificamente a EFD ICMS IPI. A motivação para o
desenvolvimento desta solução fundamentou-se na rigidez dos sistemas
tradicionais, que frequentemente exigem refatoração de código-fonte a cada nova
atualização normativa do fisco.

Os resultados obtidos ao longo das etapas de homologação confirmam que a
hipótese arquitetural proposta é tecnicamente viável e altamente vantajosa.
A adoção de um modelo de metadados, baseado em arquivos JSON, permitiu
desacoplar as regras de negócio da lógica de extração, apesar da necessidade de
extração manual dessas regras de negócios a partir dos documentos oficiais.

Adicionalmente, as técnicas de validação utilizadas, mesmo que ainda
embrionárias, atestaram a integridade do processo. A precisão verificada na
remontagem dos arquivos brutos, a partir dos dados em memória e do banco de
dados relacional, assegura que o \textit{pipeline} cumpre o rigor exigido por
sistemas de natureza fiscal, garantindo zero perda informacional. Contudo,
é válido salientar que essa ainda não é uma abordagem ideal, muito menos
escalável. A natureza \textit{ad-hoc} e manual da execução dos testes,
além de demasiadamente custosa e propensa a erros, pode ser melhorada e
explorada, buscando atender a métricas de cobertura de código, casos de teste de
regressão que permitam testar melhorias de código e, principalmente,
fornecer um modelo expansível e automatizado de execução dos mesmos.

Além disso, o desenvolvimento expôs desafios significativos inerentes à
engenharia de dados em larga escala. A identificação dos gargalos de paginação
do tipo JSONB e a degradação temporal causada pela necessidade do
\textit{PostgreSQL} de manter a consistência dos dados evidenciaram as
limitações de ferramentas relacionais puras para armazenamento de formato
bruto, bem como dos \textit{trade-offs} de modelagem. Contudo, a análise
detalhada desses comportamentos imprevisíveis constitui uma contribuição
acadêmica e técnica relevante deste trabalho, balizando o entendimento sobre
escolhas arquiteturais e fundamentando as diretrizes para a evolução contínua
da ferramenta.

Em suma, o projeto atendeu de maneira satisfatória aos requisitos propostos,
entregando um orquestrador resiliente, rastreável e flexível. A solução
consolida uma base sólida que não apenas mitiga o esforço de manutenção de
\textit{software} perante a volatilidade da legislação tributária brasileira,
mas também habilita a transformação de dados opacos em informações estruturadas
prontas para consumo analítico e a aplicação direta de uma ideia que pode ser
expandida para outros sistemas de mesmo gênero.

\subsection{Trabalhos Futuros}
\label{sec:trabalhos_futuros}

Considerando as limitações mapeadas durante a avaliação do protótipo e o vasto
potencial de expansão da arquitetura estabelecida, sugerem-se as seguintes
frentes de pesquisa e desenvolvimento para o aprimoramento da solução:

\begin{itemize}
    \item \textbf{Otimização do Motor de Persistência para Volumes Massivos:}
    Tendo em vista que a saturação de desempenho está atrelada ao custo de
    manutenção da consistência transacional do SGBD, sugere-se o desenvolvimento
    de um pipeline otimizado de carga. Este estudo futuro envolveria a adoção de
    \textit{tabelas de estagiamento}\footnote{Estruturas temporárias de
    armazenamento usadas para receber dados brutos
    rapidamente, minimizando o custo de restrições de integridade e atualização
    de índices durante a carga inicial.} sem índices para a
    recepção bruta dos registros EFD ICMS IPI via
    \textit{bulk inserts}\footnote{Técnica de ingestão de dados que envia
    grandes volumes de registros em um único lote ou comando para o SGBD,
    reduzindo o custo computacional por linha inserida.},
    seguida pela consolidação assíncrona e particionamento dos dados nas
    tabelas finais do modelo relacional;

    \item \textbf{Testes de Estresse e Escalabilidade Horizontal:} Executar
    baterias de testes quantitativos com volumetria de produção em ambiente
    distribuído, com o objetivo de traçar curvas de escalabilidade linear
    do Orquestrador, avaliando o comportamento do gerenciamento concorrente
    de \textit{workers} sob carga massiva e extrema;

    \item \textbf{Utilização de \textit{parsing} e \textit{loading} baseado
    em metadados em outros cenários:} Visando atestar a eficiência e
    manutenibilidade dessa abordagem, é sugerido que se faça o uso da
    mesma para outros formatos de documento em texto semiestruturado que
    requiram a criação de um \textit{parsing} dedicado e que tenham alta
    volatilidade. O objetivo aqui seria entender como essa abordagem pode ser
    aplicada em outros casos e quais os níveis de sucesso e mudanças necessárias
    a serem aplicadas;
\end{itemize}